% Notes-only consequence proposed after the separate D94 stratum addendum.
% Unregistered. No canonical source or owner row is changed.
\subsubsection{Exhaustion for existing locally primitive unit sections}

Let \(s\) be an \emph{actual} section of the geometrically
pro-five metabelian tripod quotient of D.35. Suppose that its
abelian pair has zero valuation vector at five and a primitive
local unit vector:
\[
 (a_5,b_5)=(0,0),\qquad
 z=(\alpha,\beta)\in\mathbb Z_5^2\setminus5\mathbb Z_5^2.
                                                        \tag{LP.1}
\]
There is no finite-support assumption. Write \(L\) for its
actual closed away-five valuation lattice. Then D.81, D.83,
and the separately proposed D.94 stratum/endpoint addendum
exhaust the compact kernels of all such sections. This is a
conditional-on-existence classification of their cohomology
modules, not a criterion asserting that arbitrary pairs lift.

\begin{enumerate}
\item If \(L\) has rank one, then
\(\mathcal S_s=\mathcal K_s=0\).
\item If \(L\bmod5\) has rank two, then \(L=\mathbb Z_5^2\).
If the residual line of \(z\) is none of the three peripheral
lines, \(\mathcal S_s=\mathcal K_s=\mathbb Z/5\).
If it is peripheral, \(\mathcal S_s=0\), and
\(\mathcal K_s=Z(d)\), where
\[
 d=v_5(\alpha)+v_5(\beta)+v_5(\alpha-\beta),
 \quad Z(d)=\mathbb Z_5^d\ (d<\infty),\quad
 Z(\infty)=\prod_{\mathbb N}\mathbb Z_5.
                                                        \tag{LP.2}
\]
Here in the peripheral case \(d\ge1\), and at most one of
the three summands is positive or infinite.
\item In the remaining case, \(L\) has rank two and residual
rank one. After an actual permutation of the three punctures,
\[
 L=5^r\mathbb Z_5\oplus\mathbb Z_5,\quad r\ge1,
 \qquad z=(\alpha,\beta),\quad\beta\in\mathbb Z_5^\times.
                                                        \tag{LP.3}
\]
Put \(N=v_5(\alpha)\ge r\). For finite \(N\),
\[
 \begin{aligned}
 \mathcal S_s&\simeq
       \bigoplus_{j=1}^{r-1}\mathbb Z/5^{N-r+j}
                  \oplus\mathbb Z/5^{N+1},\\
 \mathcal K_s&\simeq\mathcal S_s\oplus\mathbb Z_5^{N-r}.
 \end{aligned}                                         \tag{LP.4}
\]
For \(N=\infty\),
\(\mathcal S_s\simeq\mathbb Z_5^r\) and
\(\mathcal K_s\simeq\prod_{\mathbb N}\mathbb Z_5\).
These last isomorphisms are abstract compact-module
identifications; the natural kernel maps remain those in
D.94 and its endpoint proof.
\end{enumerate}

\paragraph{Proof of the exhaustion and its actual marking.}
D.87 gives \(z\in L\), with the literal uniformly convergent
local logarithm sum. Thus \(L\bmod5\ne0\). In particular
\(L\) cannot have rank zero or residual rank zero. Its
possible ranks leave exactly the three cases listed above.

Choose an away-five prime whose actual inertia vector is
nonzero modulo five; such a prime exists because otherwise
all vectors, their closed span, and their convergent local
sum would lie in \(5\mathbb Z_5^2\). Restrict the \emph{actual
metabelian section} to its local decomposition group. D.42
forces this inertia vector to lie on one of the three exact
peripheral lines, not merely on a line with the same residual
slope. Since the vector is primitive, its \(\mathbb Z_5\)
span is the entire primitive peripheral line.

If \(L\) has rank one, it therefore is that primitive line.
The primitive vector \(z\in L\) generates \(L\), so D.83
applies with equality of the actual global and local operator
images. It gives the first conclusion, including the actual
finite-level localization isomorphisms used to identify the
compact terms.

If \(L\) has rank two and residual rank one, the primitive
inertia vector just found has the unique nonzero residual
direction of \(L\). Use an actual rational automorphism of
the tripod permuting the three punctures to send its exact
peripheral line to the second coordinate axis. The induced
maps on abelian Kummer pairs are generated by swapping the
coordinates and by \((a,b)\mapsto(-a,b-a)\); the sign
Kummer term vanishes at five. As in D.79, transport the
actual tangential frame, and then change to the standing
frame by the corresponding inner gauge. This gives an actual
isomorphism of the twisted Alexander modules and their whole
global/local cohomology diagrams. It is not an assertion of
invariance under an arbitrary linear coordinate change.

Now \(\mathbb Z_5 e_2\subset L\). The projection of \(L\)
to the first coordinate is a nonzero proper closed ideal
\(5^r\mathbb Z_5\), with \(r\ge1\). Subtracting the
second coordinate using the contained axis proves exactly
\(L=5^r\mathbb Z_5e_1\oplus\mathbb Z_5e_2\).
Since \(z\) is primitive and its first coordinate is
five-divisible, its second coordinate is a unit. These are
precisely the complete hypotheses of the D.94 addendum,
which proves (LP.4) and its infinite endpoint.

Finally if \(L\bmod5\) has rank two, compact Nakayama gives
\(L=\mathbb Z_5^2\), and D.81 applies directly. The primitive
hypothesis on \(z\) reduces that theorem to the stated
peripheral/nonperipheral alternatives. No new existence,
first-cohomology rank, or higher-derived lift claim is used.
